\section{RENCANA DOKUMENTASI DAN DEPLOYMENT}

\subsection{Rencana Deployment}

Savora akan di-deploy menggunakan strategi berikut setelah development selesai:

\begin{enumerate}[leftmargin=*]
    \item \textbf{Frontend} : Deployment ke platform seperti Vercel atau Netlify dengan continuous deployment dari GitHub
    \item \textbf{Backend} : Deployment ke cloud provider seperti Railway, Render, atau AWS dengan auto-scaling capability
    \item \textbf{Database} : PostgreSQL managed instance (Supabase, Railway, atau AWS RDS)
    \item \textbf{ML Service} : Python microservice di Docker container
\end{enumerate}

\textbf{STATUS}: Deployment belum dilakukan. Aplikasi masih dalam tahap perencanaan dan akan di-deploy setelah development selesai.

\subsection{Rencana Dokumentasi Pengguna}

Dokumentasi akan dibuat untuk memudahkan pengguna menggunakan aplikasi setelah deployment:

\subsubsection{User Guide untuk Customer}
Akan mencakup:
\begin{itemize}[leftmargin=*]
    \item Cara registrasi dan login
    \item Cara mencari dan membeli produk
    \item Cara memahami Food Eligibility Score
    \item Cara tracking pesanan
    \item Cara memberikan review
    \item Cara melihat impact dashboard
\end{itemize}

\subsubsection{User Guide untuk UMKM}
Akan mencakup:
\begin{itemize}[leftmargin=*]
    \item Cara registrasi dan verifikasi
    \item Cara menambahkan produk
    \item Cara mengisi Food Eligibility Assessment
    \item Cara menggunakan Dynamic Pricing
    \item Cara mengelola pesanan
    \item Cara membaca Food Rescue Score
\end{itemize}

\subsubsection{User Guide untuk Admin}
Akan mencakup:
\begin{itemize}[leftmargin=*]
    \item Cara mengakses admin panel
    \item Cara verifikasi UMKM baru
    \item Cara monitoring platform
    \item Cara moderasi konten
    \item Cara mengelola user
\end{itemize}

\subsection{Rencana Akun Demo}

Untuk keperluan testing dan penilaian kompetisi, akan disediakan akun demo:

\begin{table}[H]
\centering
\caption{Rencana Akun Demo}
\label{tab:akun_demo}
\begin{tabularx}{\textwidth}{|l|X|X|}
\hline
\textbf{Role} & \textbf{Email} & \textbf{Purpose} \\
\hline
Admin & admin@resqfood.demo & Testing admin features \\
\hline
UMKM & umkm.demo@resqfood.demo & Testing UMKM workflow \\
\hline
Customer & customer.demo@resqfood.demo & Testing customer flow \\
\hline
\end{tabularx}
\end{table}

\textbf{Sample Data yang Akan Di-seed:}
\begin{itemize}[leftmargin=*]
    \item 10+ UMKM dengan profil lengkap
    \item 30+ produk makanan surplus dengan berbagai kategori
    \item 50+ transaksi historis untuk testing
    \item Review dan rating dari customer
    \item Statistik dashboard (food waste saved, carbon impact)
\end{itemize}

\textbf{CATATAN}: Akun demo dan sample data di atas adalah RENCANA. Belum ada aplikasi yang di-deploy saat ini.

\subsection{Requirement Sistem}

Aplikasi Savora akan memiliki requirement minimum berikut:

\begin{table}[H]
\centering
\caption{Minimum System Requirements}
\label{tab:requirements}
\begin{tabularx}{\textwidth}{|l|X|}
\hline
\textbf{Komponen} & \textbf{Spesifikasi Minimum} \\
\hline
Browser & Chrome 90+, Firefox 88+, Safari 14+, Edge 90+ \\
\hline
Koneksi Internet & 2 Mbps untuk pengalaman optimal \\
\hline
Device & Desktop, laptop, tablet, atau smartphone \\
\hline
Screen Resolution & 360px width minimum (mobile-first) \\
\hline
JavaScript & Enabled (required) \\
\hline
\end{tabularx}
\end{table}

\subsection{Rencana FAQ dan Troubleshooting}

Dokumentasi FAQ dan troubleshooting akan dibuat mencakup:

\subsubsection{Common Issues yang Akan Didokumentasikan}
\begin{itemize}[leftmargin=*]
    \item Masalah login dan authentication
    \item Masalah loading halaman
    \item Masalah upload gambar
    \item Masalah pembayaran (jika ada)
    \item Masalah notifikasi
    \item Masalah perhitungan FES
\end{itemize}

Setiap issue akan didokumentasikan dengan:
\begin{enumerate}[leftmargin=*]
    \item Deskripsi masalah
    \item Penyebab umum
    \item Langkah-langkah solusi
    \item Cara menghubungi support jika tidak terselesaikan
\end{enumerate}

\subsection{Rencana Support dan Kontak}

Setelah deployment, akan tersedia channel support:

\begin{itemize}[leftmargin=*]
    \item \textbf{Email Support} : Untuk pertanyaan dan issue teknis
    \item \textbf{GitHub Issues} : Untuk bug reports dan feature requests
    \item \textbf{Documentation Website} : Knowledge base dan tutorials
\end{itemize}

\subsection{Dokumentasi Teknis untuk Developer}

Akan dibuat dokumentasi teknis mencakup:

\subsubsection{API Documentation}
\begin{itemize}[leftmargin=*]
    \item Daftar semua endpoints dengan method dan parameters
    \item Request/response examples
    \item Authentication requirements
    \item Error codes dan handling
    \item Rate limiting policies
\end{itemize}

\subsubsection{Database Schema Documentation}
\begin{itemize}[leftmargin=*]
    \item Entity Relationship Diagram (ERD)
    \item Table descriptions
    \item Field definitions dan constraints
    \item Relationships dan foreign keys
    \item Indexes dan performance considerations
\end{itemize}

\subsubsection{Deployment Guide}
\begin{itemize}[leftmargin=*]
    \item Environment setup instructions
    \item Configuration requirements
    \item Database migration steps
    \item CI/CD pipeline setup
    \item Monitoring dan logging setup
\end{itemize}

\subsection{Rencana Maintenance Documentation}

Dokumentasi maintenance akan mencakup:

\begin{itemize}[leftmargin=*]
    \item \textbf{Backup Strategy} : Frequency, retention policy, restore procedures
    \item \textbf{Monitoring} : Key metrics to track, alerting thresholds
    \item \textbf{Scaling Guidelines} : When and how to scale infrastructure
    \item \textbf{Security Updates} : Process for patching vulnerabilities
    \item \textbf{Performance Optimization} : Common bottlenecks dan solutions
\end{itemize}

\subsection{Timeline Dokumentasi}

\begin{table}[H]
\centering
\caption{Rencana Timeline Dokumentasi}
\label{tab:doc_timeline}
\begin{tabularx}{\textwidth}{|l|X|l|}
\hline
\textbf{Fase} & \textbf{Deliverable} & \textbf{Timing} \\
\hline
Development & API documentation, code comments & Ongoing \\
\hline
Pre-deployment & User guides, admin documentation & 1-2 minggu \\
\hline
Post-deployment & FAQ, troubleshooting guide & Iteratif \\
\hline
Maintenance & Technical documentation updates & Ongoing \\
\hline
\end{tabularx}
\end{table}

\textbf{CATATAN PENTING}: Semua dokumentasi di atas adalah RENCANA yang akan dibuat setelah aplikasi selesai dibangun dan di-deploy. Saat ini aplikasi belum ada, sehingga belum ada user guide, FAQ, atau troubleshooting documentation yang dapat diberikan. Dokumentasi akan dibuat bertahap seiring dengan progress development dan feedback dari user testing.
